% part: intuitionistic-logic
% chapter: propositions-as-types
% section: types

\documentclass[../../../include/open-logic-section]{subfiles}

\begin{document}

\olfileid{int}{pty}{typ}

\olsection{પ્રકારો}

$(MN)$ જેવા સાબિતી-પદથી એકલેથી તે જે સાબિતીનું પદ છે
તે સાબિતી અનન્ય રીતે નક્કી થતી નથી.
કારણ કે મુક્ત ચલો, એટલે સ્વયંસિદ્ધોનાં સાબિતી-પદો,
અનુરૂપ સ્વયંસિદ્ધોમાંનાં !!{formula}s જણાવતા નથી.
પરંતુ તેમને સંદર્ભ~$\Gamma$માં જણાવીએ અને પદ યોગ્ય રીતે
પ્રકારિત હોય, તો સાબિતી અનન્ય રીતે નક્કી થાય છે.

\begin{defn}
સાબિતી-પદ~$N$ માટેનો \emph{સંદર્ભ} એવો સમૂહ $\Gamma$ છે
કે~$N$ના દરેક મુક્ત ચલ માટે
$\Gamma$માં બરાબર એક જ $x : !A$ની જોડી હોય.
\end{defn}

નોંધો કે વ્યાખ્યા મુજબ $\Gamma$માં
$x:!A$ની જોડી \emph{ફક્ત}~$N$ના મુક્ત ચલો માટે જ
હોવી જરૂરી નથી.
દાખલા તરીકે $\Gamma = \{x:!A, y:!B, z:!C\}$
એ $\pair{x}{y}$ માટેનો સંદર્ભ છે.

\begin{defn}\ollabel{defn:type}
  $\Gamma$ એ~$N$ માટેનો સંદર્ભ છે એમ ધારો.
  $N$નો \emph{પ્રકાર} ($\Gamma$ની સાપેક્ષે)
  આ રીતે આગમનથી વ્યાખ્યાયિત થાય છે:
  \begin{enumerate}
  \item $x$નો પ્રકાર $!A$ છે જો અને ફક્ત જો $x:!A \in \Gamma$.
  \item $N$નો પ્રકાર $x:!A, \Gamma$ની સાપેક્ષે~$!B$ હોય,
    તો $\lambd[\typeof{x}{!A}][N]$નો પ્રકાર $!A \lif !B$ છે.
  \item $N$નો પ્રકાર $!A \lif !B$ અને $M$નો પ્રકાર~$!A$ હોય,
    તો $(NM)$નો પ્રકાર~$!B$ છે.
  \item $N$નો પ્રકાર $!A$ અને $M$નો પ્રકાર~$!B$ હોય,
    તો $\pair{N}{M}$નો પ્રકાર $!A \land !B$ છે.
  \item $N$નો પ્રકાર~$!A_1 \land !A_2$ હોય,
    તો $\proj{i}{N}$નો પ્રકાર~$!A_i$ છે.
  \item $N$નો પ્રકાર $!A$ હોય, તો
    $\inj[!B]{i}{N}$નો પ્રકાર $i=1$ હોય ત્યારે $!A \lor !B$
    અને $i=2$ હોય ત્યારે $!B \lor !A$ છે.
    \sourcecorrection{OLMD-297}{મૂળમાં આ નિયમના આધારનું પદ $N$ છે પરંતુ વિકલ્પ-દાખલમાં $!M$ લખ્યું છે. દાખલનું પદ એ જ $N$ કર્યું છે, જેમ બંને અનુરૂપ અનુમાન-નિયમોમાં છે.}
  \item $M$નો પ્રકાર $!A \lor !B$ હોય,
    $N_1$નો પ્રકાર $x:!A, \Gamma$ની સાપેક્ષે $!C$
    અને $N_2$નો પ્રકાર $y:!B, \Gamma$ની સાપેક્ષે~$!C$ હોય,
    તો $\dcase{M}{x}{N_1}{y}{N_2}$નો પ્રકાર~$!C$ છે.
    \sourcecorrection{OLMD-298}{મૂળમાં શાખાઓના સંદર્ભમાં $x$ અને $y$ છે પરંતુ કેસ-પદમાં $x_1$ અને $x_2$ છે. કેસ-પદમાં શાખાઓનાં એ જ $x$ અને $y$ બદ્ધ ચલ રાખ્યાં છે.}
  \item $N$નો પ્રકાર~$\lfalse$ હોય,
    તો $\abort{!A}{N}$નો પ્રકાર~$!A$ છે.
  \end{enumerate}
\end{defn}

\begin{prop}
દરેક યોગ્ય રીતે પ્રકારિત સાબિતી-પદ~$N$ને
$N$ માટેના સંદર્ભ~$\Gamma$ની સાપેક્ષે બરાબર એક પ્રકાર હોય છે.
\sourcecorrection{OLMD-299}{મૂળ દરેક સાબિતી-પદ માટે પ્રકારનું અસ્તિત્વ તથા અનન્યતા કહે છે. પરંતુ આ જ વિભાગનું સ્વ-પ્રયોજનનું ઉદાહરણ બતાવે છે કે પદને પ્રકાર મળી ન શકે. તેથી અહીં યોગ્ય રીતે પ્રકારિત પદની શરત સ્પષ્ટ કરી છે; આગમન પ્રકારની અનન્યતા આપે છે, મનસ્વી પદ માટે પ્રકારનું અસ્તિત્વ નહીં. શરૂઆતના સાબિતીની અનન્યતા અંગેના વાક્યમાં પણ આ શરત સ્પષ્ટ કરી છે.}
\end{prop}

\begin{proof}
  $N$ પર આગમનથી.
\end{proof}

$N$નો પ્રકાર $\Gamma$ની સાપેક્ષે~$!A$ હોય,
તો $\Gamma \Sequent N : !A$ લખીએ.
\olref{defn:type}ની કલમો પરિચિત લાગવી જોઈએ:
તેઓ એક નાના તફાવત સાથે બરાબર \Log{tN2i}ના નિયમો જ છે;
સંદર્ભ~$\Gamma$ બધે એ જ છે.
આ વ્યાખ્યાને \olref{tab:tN2ip}માં આપેલા નિયમો ધરાવતા
\Log{tN3i} કલન તરીકે રજૂ કરી શકીએ.

\subfile{rules-tN3}

સાબિતી-પદો $\lambd$-કલનના એક સ્વરૂપનાં પદો છે,
જેને \emph{ગુણાકારો, સરવાળાઓ અને ખાલી પ્રકાર ધરાવતું
પ્રકારિત $\lambd$-કલન} કહે છે.
શાસ્ત્રીય $\lambd$-કલનમાં ફક્ત ચલોથી,
\emph{$\lambd$-અમૂર્ત પદો}~$\lambd[x][N]$થી
અને પ્રયોજન $(MN)$થી બનેલાં પદો હોય છે.
તેમાં પ્રકારિત $\lambd$-અમૂર્ત પદ
$\lambd[\typeof{x}{!A}][N]$ની~$!A$ની નોંધ હોતી નથી.
$\lambd$-કલન સંગણનાનું એવું નિદર્શ છે જે ટ્યુરિંગ-સંપૂર્ણ છે
(એટલે દરેક આંશિક સંગણનીય ફલનને તેમાં કોઈ પદથી રજૂ કરી શકાય).
સાબિતી-પદો અને પ્રકાર-નિયમોથી આપેલું $\lambd$-કલનનું સ્વરૂપ
સંગણનાનું મર્યાદિત નિદર્શ છે, પરંતુ મૂળભૂત વિચારો એ જ છે.
પ્રકાર-સોંપણી તેને મર્યાદિત કરે છે:
ફક્ત \emph{યોગ્ય રીતે પ્રકારિત} પદો,
એટલે સંદર્ભ~$\Gamma$ની સાપેક્ષે પ્રકાર ધરાવતાં સાબિતી-પદો,
માન્ય પદો છે.
દાખલા તરીકે શાસ્ત્રીય અપ્રકારિત $\lambd$-કલનમાં
$\lambd[x][(xx)]$ માન્ય પદ છે, પરંતુ
કોઈ પણ $!A$ માટે $\lambd[\typeof{x}{!A}][(xx)]$
યોગ્ય રીતે પ્રકારિત નથી.
કારણ કે $(xx)$ને પ્રકાર મળવા માટે પ્રકાર-નિયમો મુજબ
પહેલા~$x$નો પ્રકાર $!A \lif !B$
અને બીજા $x$નો પ્રકાર~$!A$ હોવો જોઈએ.
આ અશક્ય છે, કારણ કે $!A$ એ~$!A \lif !B$નું
યથાર્થ ઉપ-!!{formula} છે.

સહજ રીતે પ્રકારને તે પ્રકારના પદથી નામિત થતી વસ્તુનો
પ્રકાર ગણી શકીએ.
એટલે પ્રકાર $!A \land !B$નું પદ એવી જોડી પસંદ કરે છે
જેનો પહેલો ઘટક પ્રકાર~$!A$નો અને બીજો ઘટક પ્રકાર~$!B$નો હોય,
એટલે ગુણાકાર~$!A \times !B$નું !!a{element}.
પ્રકાર $!A \lif !B$નું પદ પ્રકાર $!A$ની વસ્તુઓમાંથી
પ્રકાર~$!B$ની વસ્તુઓમાં જતું ફલન છે.
પ્રકાર $!A \lor !B$નું પદ પ્રકાર $!A$ની કે
પ્રકાર~$!B$ની વસ્તુ છે, સાથે તે બેમાંથી કઈ છે તેની માહિતી હોય છે.
આ બે સમૂહો $!A$ અને $!B$ના વિયુક્ત સંઘ જેવું છે,
જેને $\{ \tuple{i,x}
: i = 1 \text{ અથવા } 2, x \in !A \text{ જો } i = 1, x \in !B \text{ જો
} i=2\}$ તરીકે વ્યાખ્યાયિત કરીએ.
પ્રકાર $\lfalse$ ખાલી સમૂહ છે, એટલે પદનો એ પ્રકાર હોય,
તો તે કોઈ વસ્તુને નામિત કરી શકતું નથી.
સ્વાભાવિક નિગમન (\Log{tN3i} સહિત) સુસંગત હોવાથી,
ખુલ્લી ધારણાઓ ન હોય ત્યાં સુધી~$\lfalse$ની !!{proof} નથી.
એટલે પ્રકાર~$\lfalse$નું કોઈ બંધ સાબિતી-પદ
(મુક્ત ચલ વિનાનું પદ) નથી.

પ્રકારિત $\lambd$-કલનના સંકારકો સહજ રીતે
યોગ્ય પ્રકારની વસ્તુને નામિત કરતાં પદો આપે છે,
જો તેમના પર લાગુ પાડવામાં આવતાં પદો અમુક વસ્તુઓને નામિત કરતાં હોય.
દાખલા તરીકે ``$N_1 : !A$ અને $N_2 : !B$ હોય,
તો $\pair{N_1}{N_2} : !A \land !B$'' એવું કહે છે
કે $N_1$ પ્રકાર~$!A$ની વસ્તુને અને
$N_2$ પ્રકાર~$!B$ની વસ્તુને નામિત કરે,
તો $\pair{N_1}{N_2}$ પ્રકાર $!A \land !B$ની
વસ્તુને નામિત કરે છે.

\end{document}
